\documentclass[11pt]{article}
\usepackage[T1]{fontenc}
\usepackage[a4paper,margin=24mm]{geometry}
\usepackage{amsmath,amssymb}
\usepackage[hidelinks]{hyperref}
\usepackage{microtype}

\setlength{\parindent}{0pt}
\setlength{\parskip}{0.58em}

\title{Second Quantisation for General Infinitely Divisible Measures:\\
A Scoped Later Answer}
\author{Literature status for arXiv:1305.5146}
\date{13 August 2026}

\begin{document}
\maketitle

\begin{center}
\fbox{\parbox{0.91\textwidth}{
\textbf{Status: literature already answered, with a scope correction.}
Applebaum--van Neerven's follow-up arXiv:1405.1276 explicitly completes the
Gaussian/pure-jump programme for general infinitely divisible measures when
the skew factor is infinitely divisible.  The same follow-up proves that this
factor hypothesis is not automatic.
}}
\end{center}

\section*{Source conjecture}

For a skew map $T:E_1\to E_2$ with
\[
 T(\mu_1)*\rho=\mu_2,
\]
the source constructs a contraction
\[
 (P_Tf)(x)=\int_{E_2}f(Tx+y)\,d\rho(y)
\]
and identifies it with symmetric second quantisation when the endpoint laws
are both Gaussian (Theorem~3.5) or both infinitely divisible of pure-jump
type (Theorem~4.4)~\cite{source}.

Remark~4.5 on source PDF page~14 asks whether these theorems extend when
$\mu_1$ and $\mu_2$ are arbitrary infinitely divisible measures.

\section*{Exact later answer}

The same authors' separate follow-up explicitly cites the source and states on
supporting PDF page~2:
\begin{quote}
In this article, we complete the programme by dealing with the case where the
$\lambda_i$ are general infinitely divisible measures.
\end{quote}
It decomposes each law canonically as
\[
 \lambda_i=\gamma_i*\pi_i,
\]
where $\gamma_i$ is centred Gaussian and $\pi_i$ is the drift-plus-pure-jump
part.  Proposition~2.1 shows that, if the skew factor is infinitely divisible,
the skew relation splits into Gaussian and pure-jump skew relations.  The two
known second-quantisation maps can therefore be tensored.

Theorem~3.7 on supporting PDF page~10 gives the resulting commuting diagram.
Its one-particle Hilbert space is
\[
 H_{\gamma_i}\oplus L^2(E_i,\nu_i),
\]
and its Fock-space arrow is
\[
 \Gamma(T^*)=
 \bigoplus_{n=0}^{\infty}(T^*)^{\odot n}.
\]
The transition operator $P_T$ is intertwined with this arrow after embedding
$f$ as $F_f(x,y)=f(x+y)$ into the independent Gaussian--jump product space.
This is precisely the mixed extension suggested in the source~\cite{answer}.

\section*{Why the scope qualifier matters}

Theorem~3.7 uses the standing assumption that the skew factor $\rho$ is
infinitely divisible.  Endpoint infinite divisibility alone does not imply
this.  Example~2.2 (Rosi\'nski) on supporting PDF page~4 constructs an
infinitely divisible law $\mu_1$, a non-infinitely-divisible law $\rho$, and an
infinitely divisible law $\mu_2=\mu_1*\rho$.

Accordingly, the later paper gives an explicit affirmative answer to the
natural infinitely-divisible-factor formulation and identifies a genuine
obstruction to deleting that qualification.  It does not assert that the
unrestricted arbitrary-factor wording follows merely because $\mu_1$ and
$\mu_2$ are infinitely divisible.  This status note therefore records a
scoped literature answer rather than overstating Theorem~3.7.

\section*{Verification and provenance}

The source conjecture was checked on PDF page~14.  The follow-up's explicit
programme statement was checked on PDF page~2, Proposition~2.1 and the
Rosi\'nski example on pages~3--4, and the final mixed theorem on PDF page~10.
The follow-up is by the same authors, cites the source as its previous paper,
and explicitly describes its result as completing that programme.  The
relation is therefore author-explicit, not an agent-only theorem match.

\begin{thebibliography}{9}
\bibitem{source}
D. Applebaum and J. van Neerven,
\emph{Second quantisation for skew convolution products of measures in
Banach spaces},
Electron. J. Probab. \textbf{19} (2014), paper 11, 1--17;
arXiv:1305.5146.

\bibitem{answer}
D. Applebaum and J. van Neerven,
\emph{Second quantisation for skew convolution products of infinitely
divisible measures},
Infin. Dimens. Anal. Quantum Probab. Relat. Top. \textbf{18} (2015),
1550003;
arXiv:1405.1276.
\end{thebibliography}

\end{document}

